\documentclass[11pt,a4paper]{report}

% ================= PACKAGES =================
\usepackage[margin=2.5cm]{geometry}
\usepackage{graphicx}
\usepackage{xcolor}
\usepackage{titlesec}
\usepackage{fancyhdr}
\usepackage{booktabs}
\usepackage{array}
\usepackage{enumitem}
\usepackage{tcolorbox}
\usepackage{hyperref}
\usepackage{caption}
\usepackage{fontspec}
\usepackage{tocloft}
\usepackage{longtable}
\usepackage{tikz}
\usetikzlibrary{arrows.meta,positioning}

\setmainfont{Latin Modern Sans}[Scale=1.0]
\setmonofont{Latin Modern Mono}[Scale=0.9]

% ================= COLORS =================
\definecolor{primary}{HTML}{3C3489}
\definecolor{primarydark}{HTML}{26215C}
\definecolor{accent}{HTML}{1D9E75}
\definecolor{warn}{HTML}{B75E09}
\definecolor{lightgray}{HTML}{F1EFE8}
\definecolor{warnbg}{HTML}{FBF0E2}
\definecolor{bodytext}{HTML}{2C2C2A}

\color{bodytext}

% ================= SECTION STYLING =================
\titleformat{\chapter}[display]
  {\normalfont\Huge\bfseries\color{primary}}
  {\chaptertitlename\ \thechapter}{10pt}{\Huge}
\titlespacing*{\chapter}{0pt}{0pt}{30pt}

\titleformat{\section}
  {\normalfont\Large\bfseries\color{primarydark}}{\thesection}{1em}{}
\titleformat{\subsection}
  {\normalfont\large\bfseries\color{primary}}{\thesubsection}{1em}{}

% ================= HEADER / FOOTER =================
\pagestyle{fancy}
\fancyhf{}
\renewcommand{\headrulewidth}{0.4pt}
\fancyhead[L]{\small\textcolor{primarydark}{SmartSchool -- School Management ERP}}
\fancyhead[R]{\small\textcolor{primarydark}{Software Requirements Specification}}
\fancyfoot[C]{\thepage}

% ================= BOXES =================
\newtcolorbox{infobox}[1][]{
  colback=lightgray, colframe=primary, boxrule=0.6pt, arc=2pt,
  left=8pt, right=8pt, top=6pt, bottom=6pt, #1
}
\newtcolorbox{futurebox}[1][]{
  colback=warnbg, colframe=warn, boxrule=0.6pt, arc=2pt,
  left=8pt, right=8pt, top=6pt, bottom=6pt, #1
}

% ================= HYPERREF =================
\hypersetup{
  colorlinks=true,
  linkcolor=primary,
  citecolor=primary,
  urlcolor=accent,
  pdftitle={SmartSchool SRS - SEN4121 Large System Environment},
  pdfauthor={Group 2 -- ICT University Yaounde}
}

\captionsetup{font=small,labelfont={bf,color=primary}}

\newcommand{\sysname}{SmartSchool}
\newcolumntype{L}[1]{>{\raggedright\arraybackslash}p{#1}}

\begin{document}

% ================= TITLE PAGE =================
\begin{titlepage}
\begin{center}
\vspace*{1.2cm}

{\color{primary}\rule{\linewidth}{2pt}}

\vspace{0.8cm}
{\Huge\bfseries\color{primary} Software Requirements Specification}

\vspace{0.4cm}
{\LARGE\color{primarydark} \sysname{} --- Student Records, Fee Invoicing\\[2pt] \& Staff Payroll Management System}

\vspace{0.6cm}
{\color{primary}\rule{\linewidth}{2pt}}

\vspace{2cm}

{\Large Prepared for: \textbf{ICT University Ya\'ound\'e}}

\vspace{0.3cm}
{\large Course: \textbf{SEN4121 -- Large System Environment} \quad (Practical Examination, Summer 2026)}

\vspace{0.3cm}
{\large Examiner: Engr. Tanwi Nkiamboh}

\vspace{1.5cm}

\begin{tabular}{ll}
\textbf{Document version:} & 1.0 \\[4pt]
\textbf{Group number:} & Group 2 \\[4pt]
\textbf{Date:} & \today \\[4pt]
\textbf{Status:} & Living document -- updated through Week 4 \\
\end{tabular}

\vspace{1cm}

\begin{tabular}{@{}p{6cm}p{4.5cm}@{}}
\toprule
\textbf{Group Member} & \textbf{Matricule} \\
\midrule
Kwete Rayan & ICTU20241377 \\
Kamdeu Yamdjeuson Neil Marshall & ICTU20241386 \\
Younga Tchappi Dylane & ICTU2241836 \\
Fru Chi Ehud Neba & ICTU20241812 \\
\bottomrule
\end{tabular}

\vfill

{\small\color{bodytext} This document specifies the functional and non-functional
requirements for \sysname{}, a microservice-based ERP for a school, covering
student academic records, fee invoicing, and staff payroll, built as part of
the SEN4121 6-week practical examination series.}

\end{center}
\end{titlepage}

\tableofcontents
\clearpage

% ================================================================
\chapter{Introduction}

\section{Purpose}
This Software Requirements Specification (SRS) defines the functional and
non-functional requirements for \sysname{}, a school Enterprise Resource
Planning (ERP) system built as a set of cooperating microservices behind a
single API Gateway. It restates the system's required behaviour as
verifiable requirements, independent of implementation detail, and is
intended as an agreed basis between the development group and the course
examiner for what the system must do across the six weekly practical
sessions of SEN4121.

Unlike a purely aspirational specification, every requirement below is
tagged with a \textbf{status}: \texttt{Implemented} (already built and
demonstrable) or \texttt{Planned} (required by a later week of the exam
series but not yet built). This keeps the document honest about the current
state of the system rather than describing a system that does not yet
exist.

\section{Scope}
\sysname{} covers the three functional areas required by the exam brief:
\begin{itemize}[itemsep=2pt]
    \item \textbf{Academic} -- student records, courses, enrollments, and grades.
    \item \textbf{Finance} -- fee invoicing, automatically triggered by enrollment.
    \item \textbf{HR / Payroll} -- staff records and salary payslips.
\end{itemize}
The system serves three internal actors -- \textbf{Admin}, \textbf{Teacher},
and \textbf{Student} -- through a React single-page application, all traffic
passing through a single API Gateway. There is currently no external
payment-gateway integration: invoices carry a mobile-money reference field
for future reconciliation, but automated payment collection is out of scope
for the current iteration (see Section~\ref{sec:constraints}).

\section{Definitions, Acronyms and Abbreviations}
\begin{longtable}{L{3.3cm} L{10.7cm}}
\toprule
\textbf{Term} & \textbf{Definition} \\
\midrule
\endhead
Microservice & A small, independently deployable service responsible for one job (e.g. \texttt{auth-service} only handles identity). \\
API Gateway & The single entry point (\texttt{gateway}) that receives every external request and forwards it to the correct backend service. \\
RBAC & Role-Based Access Control -- what a user may do depends on their role (\texttt{admin}, \texttt{teacher}, \texttt{student}). \\
JWT & JSON Web Token -- a signed token issued at login that proves identity on later requests without a database lookup. \\
ERD & Entity-Relationship Diagram -- shows database tables, columns, and how they connect. \\
3NF & Third Normal Form -- a database design rule: no needless duplication, every column depends only on its table's primary key. \\
Message Broker & Middleman software (RabbitMQ) that lets one service publish an event another service consumes later, without a direct call. \\
Asynchronous Workflow & A process where the publishing service does not wait for the consuming service to finish. \\
OWASP Top 10 & The ten most common web application security risks. \\
Load Testing & Sending many concurrent requests to observe system behaviour under pressure. \\
CI/CD & Continuous Integration / Continuous Deployment -- automated build and test on every push. \\
SRS & Software Requirements Specification (this document). \\
SDD & Software Design Document -- the companion document describing architecture and diagrams. \\
FR / NFR & Functional Requirement / Non-Functional Requirement. \\
\bottomrule
\end{longtable}

\section{References}
\begin{itemize}[itemsep=2pt]
    \item \sysname{} Software Design Document, Version 1.0, Group 2, SEN4121, ICT University Ya\'ound\'e.
    \item SEN4121 Large System Environment -- Practical Examination brief, Summer 2026, Engr. Tanwi Nkiamboh.
    \item IEEE Std 830-1998, IEEE Recommended Practice for Software Requirements Specifications (structural reference for this document).
    \item Project source repositories: \texttt{MarketFlow} (backend services, database, documentation) and \texttt{marketflow\_frontend} (React SPA).
\end{itemize}

\section{Document Overview}
Chapter~2 gives an overall description of the product, its actors, and
constraints. Chapter~3 presents specific requirements: external interfaces
and functional requirements grouped by module, each tagged with
implementation status. Chapter~4 presents non-functional requirements.
Chapter~5 lists data requirements derived from the relational schema.
Chapter~6 traces requirements to the SDD and to the exam's weekly mark
scheme.

% ================================================================
\chapter{Overall Description}

\section{Product Perspective}
\sysname{} is a new, self-contained school-management platform composed of
cooperating services, all reached through one API Gateway:

\begin{center}
\begin{tabular}{@{}p{3.4cm}p{10.6cm}@{}}
\toprule
\textbf{Component} & \textbf{Responsibility} \\
\midrule
Gateway & Single public entry point; JWT verification, routing by URL prefix, rate limiting, combined Swagger UI. \\
Auth Service & Registration, login, JWT issuance, and the shared staff/user identity store. \\
Academic Service & Courses, enrollments, and grades; publishes an event when a student enrolls. \\
Finance Service & Invoices; consumes the enrollment event and creates an invoice asynchronously. \\
MySQL Database & Shared relational store for all modules (\texttt{smartschool} schema). \\
RabbitMQ & Message broker carrying the \texttt{enrollment.created} event between Academic and Finance. \\
React Frontend & Single-page application for registration, login, and role-scoped dashboards. \\
\bottomrule
\end{tabular}
\end{center}

Only the Gateway is reachable from outside the Docker network; every backend
service is internal-only, reflecting the "API Gateway as the only front
door" principle required by the exam brief.

\section{Product Functions}
At a high level, \sysname{} shall allow the system to:
\begin{enumerate}[itemsep=2pt]
    \item Register and authenticate Admins, Teachers, and Students, issuing a role-carrying JWT on login.
    \item Let an Admin manage courses and view system-wide reports.
    \item Let a Student browse courses, enroll, view their own enrollments, grades, and invoices.
    \item Let a Teacher record grades for the courses they teach and view those grades.
    \item Automatically generate a fee invoice whenever a student enrolls in a course, without the enrollment request waiting on it.
    \item Enforce role-based access at both the Gateway (``is this a valid, logged-in user?'') and each service (``is this specific action allowed for this role/owner?'').
    \item Protect the API from abuse via per-IP rate limiting.
    \item Compute and record staff payroll (\textit{planned, see Section~\ref{sec:constraints}}).
\end{enumerate}

\section{User Classes and Characteristics}
\begin{longtable}{L{2.6cm} L{11.4cm}}
\toprule
\textbf{User Class} & \textbf{Characteristics} \\
\midrule
\endhead
Admin & Highest-privilege internal user. Creates courses, enrolls students on their behalf, views reports, and (planned) runs payroll. Assumed to have elevated authorization and periodic use. \\
Teacher & Records grades only for enrollments in courses they teach; can view system reports alongside Admin. Assumed to be a member of staff with routine, course-scoped use. \\
Student & Registers, logs in, browses/enrolls in courses, and views their own enrollments, grades, and invoices. Assumed to have basic digital literacy. \\
\bottomrule
\end{longtable}

\section{Operating Environment}
\sysname{} is a distributed, container-based system: every backend service,
the MySQL database, and RabbitMQ run as Docker containers orchestrated by a
single \texttt{docker-compose.yml}, brought up with one command
(\texttt{docker compose up --build}). The React frontend is a separate
containerized application that talks to the Gateway over HTTP from the
browser. No specific host operating system is required beyond a working
Docker Engine.

\section{Design and Implementation Constraints}
\label{sec:constraints}
\begin{itemize}[itemsep=2pt]
    \item All inter-service authentication relies on a single shared JWT signing secret (\texttt{JWT\_SECRET}), distributed to every service via environment variable -- acceptable for a course project, not for a production multi-team deployment.
    \item The Gateway performs authentication (``is this a valid token?'') only; authorization (role/ownership checks) is deliberately left to whichever service owns the resource.
    \item Stock/academic domain rule: a course fee is captured on the \texttt{courses} table and \emph{snapshotted} onto the generated invoice at enrollment time, so a later fee change does not retroactively alter an already-issued invoice.
    \item \textbf{Payment collection is out of scope for the current iteration.} An invoice's \texttt{momo\_reference} and \texttt{method} columns exist in the schema for a future mobile-money reconciliation flow, but no external payment gateway is called by the system today; invoices are created with status \texttt{pending} and there is currently no endpoint to mark one \texttt{paid}.
    \item \textbf{Payroll is schema-only.} The \texttt{payroll\_runs} and \texttt{payslips} tables exist in the database design (Week 2) to satisfy the HR module's data model, but no service currently exposes payroll endpoints; this is planned as a Week 5 module feature (see Chapter~3, HR requirements).
\end{itemize}

\section{Assumptions and Dependencies}
\begin{itemize}[itemsep=2pt]
    \item Docker Engine and Docker Compose are available on the host running the system.
    \item RabbitMQ and MySQL containers reach a healthy state before dependent services start (enforced via Compose health checks).
    \item The frontend and Gateway are reachable from the same browser origin context (or CORS is permissive, see NFR-SEC-3).
    \item Course fee amounts and role assignments are configured by an Admin before students enroll.
\end{itemize}

% ================================================================
\chapter{Specific Requirements}

\section{External Interface Requirements}

\subsection{User Interfaces}
\begin{itemize}[itemsep=2pt]
    \item A registration/login screen for all three roles.
    \item A role-aware dashboard: every authenticated user sees the same shell, but the actions available (e.g. creating a course, recording a grade) are gated by role both in the UI and, authoritatively, by the backend.
    \item Live demo panels on the dashboard that call \texttt{/reports}, \texttt{/courses}, \texttt{/enrollments}, and \texttt{/finance/invoices} through the Gateway and render the raw JSON response, so the request/response cycle is directly observable during the oral evaluation.
\end{itemize}

\subsection{Hardware Interfaces}
None. \sysname{} is a purely software, browser/server-based system; no
specialized hardware (POS terminals, barcode scanners, biometric devices) is
required.

\subsection{Software Interfaces}
\begin{longtable}{L{3.6cm} L{10.4cm}}
\toprule
\textbf{Interface} & \textbf{Description} \\
\midrule
\endhead
MySQL 8 & Relational store for all modules; accessed by each service via a pooled \texttt{mysql2} connection using parameterized queries. \\
RabbitMQ 3.13 (AMQP 0-9-1) & Direct exchange \texttt{school\_events}; carries the \texttt{enrollment.created} event from Academic Service to Finance Service. \\
Docker / Docker Compose & Packaging and orchestration for every backend service plus the database and broker. \\
OpenAPI 3 / Swagger UI & Combined API documentation for every Gateway-routed endpoint, served at \texttt{GET /docs}. \\
\bottomrule
\end{longtable}

\subsection{Communications Interfaces}
All external communication (frontend, curl, Postman) to the system happens
over HTTP/REST with JSON bodies, through the Gateway, authenticated with a
\texttt{Bearer} JWT in the \texttt{Authorization} header where required.
Internal communication between Academic Service and Finance Service for the
enrollment/invoice workflow happens asynchronously over AMQP via RabbitMQ,
not direct HTTP.

\section{Functional Requirements}

Each requirement below is tagged \textbf{[Implemented]} (built and
demonstrable today) or \textbf{[Planned]} (required by a later exam week,
not yet built).

\subsection{Account, Authentication and Access Control}
\begin{longtable}{p{1.9cm} L{10.3cm} p{2.1cm}}
\toprule
\textbf{ID} & \textbf{Requirement} & \textbf{Status} \\
\midrule
\endhead
FR-AUTH-1 & The system shall allow a new user to register with a name, email, and password, and an optional role (\texttt{admin}, \texttt{teacher}, or \texttt{student}, defaulting to \texttt{student}). & Implemented \\
FR-AUTH-2 & The system shall store the password only as a bcrypt hash, never in plain text. & Implemented \\
FR-AUTH-3 & The system shall allow a registered user to log in with email and password and, on success, issue a JWT carrying the user's id, email, and role, expiring after 1 hour. & Implemented \\
FR-AUTH-4 & The system shall allow a client to retrieve the identity encoded in its own token via \texttt{GET /api/v1/auth/me}. & Implemented \\
FR-AUTH-5 & The system shall never include the password hash in any API response. & Implemented \\
FR-AUTH-6 & The Gateway shall verify the signature and expiry of every protected request's JWT before forwarding it to a backend service; an invalid or missing token shall be rejected with \texttt{401} at the Gateway, without reaching a backend service. & Implemented \\
FR-AUTH-7 & Each backend service shall independently re-verify the JWT and enforce role/ownership rules specific to the resource it owns (defense in depth), rather than trusting the Gateway's check alone. & Implemented \\
\bottomrule
\end{longtable}

\subsection{Academic Module -- Courses}
\begin{longtable}{p{1.9cm} L{10.3cm} p{2.1cm}}
\toprule
\textbf{ID} & \textbf{Requirement} & \textbf{Status} \\
\midrule
\endhead
FR-ACAD-1 & The system shall allow any authenticated user to list all courses. & Implemented \\
FR-ACAD-2 & The system shall allow any authenticated user to retrieve a single course by id. & Implemented \\
FR-ACAD-3 & The system shall allow only an Admin to create a new course, including its code, name, assigned teacher, credits, fee amount, and term. & Implemented \\
\bottomrule
\end{longtable}

\subsection{Academic Module -- Enrollments}
\begin{longtable}{p{1.9cm} L{10.3cm} p{2.1cm}}
\toprule
\textbf{ID} & \textbf{Requirement} & \textbf{Status} \\
\midrule
\endhead
FR-ACAD-4 & The system shall allow a Student to enroll themselves in a course. & Implemented \\
FR-ACAD-5 & The system shall allow an Admin to enroll a specified student in a course on their behalf. & Implemented \\
FR-ACAD-6 & The system shall reject a duplicate enrollment (the same student in the same course) with a clear conflict response, rather than creating a second row. & Implemented \\
FR-ACAD-7 & The system shall restrict a Student's view of enrollments to their own records only; Admin and Teacher may view or filter by student or course. & Implemented \\
FR-ACAD-8 & On a successful enrollment, the system shall publish an \texttt{enrollment.created} event to the message broker \emph{after} the enrollment is durably committed to the database, without the enrollment response waiting on the event being consumed. & Implemented \\
\bottomrule
\end{longtable}

\subsection{Academic Module -- Grades}
\begin{longtable}{p{1.9cm} L{10.3cm} p{2.1cm}}
\toprule
\textbf{ID} & \textbf{Requirement} & \textbf{Status} \\
\midrule
\endhead
FR-ACAD-9 & The system shall allow an Admin or Teacher to record a grade (quiz, assignment, midterm, or final) for a given enrollment, with a score between 0 and 100. & Implemented \\
FR-ACAD-10 & The system shall restrict a Teacher to recording grades only for enrollments in courses they themselves teach, rejecting other attempts with \texttt{403}. & Implemented \\
FR-ACAD-11 & The system shall restrict a Student's view of grades to their own; a Teacher shall see only grades for courses they teach; an Admin shall see all grades. & Implemented \\
\bottomrule
\end{longtable}

\subsection{Finance Module -- Invoicing}
\begin{longtable}{p{1.9cm} L{10.3cm} p{2.1cm}}
\toprule
\textbf{ID} & \textbf{Requirement} & \textbf{Status} \\
\midrule
\endhead
FR-FIN-1 & Upon consuming an \texttt{enrollment.created} event, the system shall create exactly one invoice for that enrollment, using the course's fee amount at the time of enrollment as a price snapshot. & Implemented \\
FR-FIN-2 & If the same enrollment event is delivered more than once (at-least-once broker delivery), the system shall not create a duplicate invoice; the duplicate is acknowledged as a successful no-op. & Implemented \\
FR-FIN-3 & The system shall allow a Student to list only their own invoices, and an Admin/Teacher to list all invoices or filter by student. & Implemented \\
FR-FIN-4 & The system shall record a mobile-money/transaction reference field on each invoice for future reconciliation. & Implemented (field only) \\
FR-FIN-5 & The system shall provide a way to mark an invoice as paid and reconcile it against a mobile-money or online payment gateway. & Planned (Wk 5/6) \\
\bottomrule
\end{longtable}

\subsection{HR Module -- Staff and Payroll}
\begin{longtable}{p{1.9cm} L{10.3cm} p{2.1cm}}
\toprule
\textbf{ID} & \textbf{Requirement} & \textbf{Status} \\
\midrule
\endhead
FR-HR-1 & The database shall maintain a payroll run record (period start/end, status) and one payslip per staff member per run (base salary, deductions, net pay). & Implemented (schema only) \\
FR-HR-2 & The system shall allow an Admin to open a payroll run for a given period and compute each staff member's net pay as gross pay minus deductions. & Planned (Wk 5) \\
FR-HR-3 & The system shall send computed payouts to staff via a payment channel and record the resulting reference on each payslip. & Planned (Wk 5/6) \\
FR-HR-4 & The system shall confirm payroll completion to the Admin with a summary once processing finishes. & Planned (Wk 5) \\
\bottomrule
\end{longtable}

\subsection{Gateway and Cross-Cutting Requirements}
\begin{longtable}{p{1.9cm} L{10.3cm} p{2.1cm}}
\toprule
\textbf{ID} & \textbf{Requirement} & \textbf{Status} \\
\midrule
\endhead
FR-GW-1 & The system shall expose exactly one publicly reachable entry point (the Gateway); no backend service shall be directly reachable from outside the container network. & Implemented \\
FR-GW-2 & The Gateway shall route each external request to the correct backend service based on its URL prefix (e.g. \texttt{/api/v1/courses} $\rightarrow$ Academic Service). & Implemented \\
FR-GW-3 & The Gateway shall limit each client IP to a configurable maximum number of requests per time window (default: 10 requests / 60 seconds) across all routes except \texttt{/health}, responding \texttt{429} once exceeded. & Implemented \\
FR-GW-4 & The system shall publish machine-readable API documentation (OpenAPI/Swagger) for every Gateway-routed endpoint, viewable at \texttt{/docs}. & Implemented \\
\bottomrule
\end{longtable}

\subsection{Reporting}
\begin{longtable}{p{1.9cm} L{10.3cm} p{2.1cm}}
\toprule
\textbf{ID} & \textbf{Requirement} & \textbf{Status} \\
\midrule
\endhead
FR-RPT-1 & The system shall allow an Admin or Teacher to view an academic performance report; a Student shall be denied access to this report with \texttt{403}. & Implemented \\
FR-RPT-2 & The role check on the reporting endpoint shall be changeable (e.g. Admin-only) without changing the authentication mechanism, and shall take effect on service restart. & Implemented \\
\bottomrule
\end{longtable}

\section{Use-Case-Driven Requirement: Enroll in a Course}
This requirement is elaborated separately because it drives the system's
central asynchronous, event-based workflow and is the most architecturally
significant use case in \sysname{}.

\begin{longtable}{L{3.3cm} L{10.7cm}}
\toprule
\textbf{Field} & \textbf{Description} \\
\midrule
\endhead
Primary actor & Student (or Admin, enrolling on a student's behalf) \\
Preconditions & Actor is authenticated with a valid JWT; the target course exists. \\
Main flow & (1) Actor submits \texttt{POST /api/v1/enrollments} with a course id through the Gateway; (2) the Gateway verifies the JWT and forwards the request to Academic Service; (3) Academic Service creates the enrollment row and commits it to MySQL; (4) Academic Service publishes an \texttt{enrollment.created} event to RabbitMQ and immediately returns \texttt{201} to the actor, without waiting on step 5; (5) Finance Service, independently and at its own pace, consumes the event and creates a pending invoice snapshotting the course's current fee. \\
Alternate flow & If the student is already enrolled in the course, the system returns \texttt{409 Conflict} and no event is published. If RabbitMQ or Finance Service is temporarily unavailable, the enrollment still succeeds; the event remains durably queued and the invoice is created once the consumer reconnects. \\
Postconditions & The enrollment exists and is visible to the student; an invoice for that enrollment eventually appears in the student's ``My invoices'' view, with no further action required from the actor. \\
\bottomrule
\end{longtable}

% ================================================================
\chapter{Non-Functional Requirements}

\section{Security}
\begin{itemize}[itemsep=3pt]
    \item[NFR-SEC-1] Every database query, in every service, shall use parameterized statements; no request-supplied value shall ever be concatenated directly into a SQL string (mitigates OWASP A03: Injection).
    \item[NFR-SEC-2] Passwords shall be hashed with bcrypt (minimum 10 salt rounds) before storage; JWTs shall carry a signed expiry of no more than 1 hour (mitigates OWASP A07: Identification and Authentication Failures).
    \item[NFR-SEC-3] Password hashes shall never appear in any API response (mitigates OWASP A02/A04: Cryptographic Failures / Insecure Design around sensitive data exposure). CORS is currently permissive across services; this is an accepted, documented gap because the system is a Bearer-token API with no ambient cookie credential to abuse.
    \item[NFR-SEC-4] The Gateway's rate limiter shall apply to authentication routes as well as data routes, so it doubles as brute-force login throttling without dedicated code.
    \item[NFR-SEC-5] Role checks shall be enforced at the service that owns the resource, not only at the Gateway, so that a compromised or bypassed Gateway check cannot alone grant unauthorized access.
\end{itemize}

\section{Performance and Scalability}
\begin{itemize}[itemsep=3pt]
    \item[NFR-PERF-1] Under a load test of at least 20 concurrent virtual users against the Gateway, the 95th-percentile response time for successful (non-rate-limited) requests shall remain under 500\,ms.
    \item[NFR-PERF-2] The rate limit window and maximum shall be externally configurable (\texttt{RATE\_LIMIT\_WINDOW\_MS}, \texttt{RATE\_LIMIT\_MAX}) without a code change, to support both a ``safety'' profile (default limits, demonstrates throttling) and a ``speed'' profile (raised limits, measures genuine backend latency) during load testing.
    \item[NFR-PERF-3] Academic Service and Finance Service shall be independently deployable and scalable processes, so that load on one module does not require scaling the other.
\end{itemize}

\section{Reliability and Consistency}
\begin{itemize}[itemsep=3pt]
    \item[NFR-REL-1] The enrollment-to-invoice workflow shall tolerate a temporarily unavailable message broker or Finance Service without failing or blocking the enrollment request (see FR-ACAD-8).
    \item[NFR-REL-2] Event messages shall be published and consumed as durable/persistent so that a broker restart does not silently lose an in-flight enrollment event.
    \item[NFR-REL-3] The event consumer shall be idempotent with respect to redelivery: processing the same event twice shall not create a duplicate invoice (see FR-FIN-2).
    \item[NFR-REL-4] Each service shall retry its connection to MySQL/RabbitMQ with backoff on startup, so that container start-order races do not cause a permanent failure.
\end{itemize}

\section{Maintainability and Portability}
\begin{itemize}[itemsep=3pt]
    \item[NFR-MAINT-1] The entire backend (database, broker, and all services) shall start from a single command (\texttt{docker compose up --build}) against a clean environment.
    \item[NFR-MAINT-2] Configuration (database credentials, JWT secret, service URLs, rate-limit parameters) shall be supplied via environment variables, never hard-coded.
    \item[NFR-MAINT-3] Each service shall own its own Dockerfile and be independently buildable and runnable in isolation for local development.
\end{itemize}

\section{Usability}
\begin{itemize}[itemsep=3pt]
    \item[NFR-USE-1] Every Gateway-routed endpoint shall be documented in a single, interactive Swagger UI reachable at one URL.
    \item[NFR-USE-2] The frontend shall clearly display the signed-in user's role and shall reflect role-based restrictions in what actions are offered, in addition to the backend's authoritative enforcement.
\end{itemize}

\section{Auditability}
\begin{itemize}[itemsep=3pt]
    \item[NFR-AUD-1] Every published domain event shall carry a unique event id, an event type, and an ISO-8601 occurrence timestamp, to support later tracing of what happened and when.
    \item[NFR-AUD-2] Every invoice shall be traceable to exactly one enrollment via a unique foreign key, preventing an enrollment from ever producing more than one invoice.
\end{itemize}

% ================================================================
\chapter{Data Requirements}

\section{Core Domain Entities}
The following entities, drawn directly from the system's relational schema
(\texttt{database/schema.sql}), define the minimum data \sysname{} must
persist to satisfy the functional requirements in Chapter~3. Full
column-level detail, the ER diagram, and the normalization rationale are
presented in the companion SDD.

\begin{longtable}{L{3.0cm} L{11.0cm}}
\toprule
\textbf{Entity} & \textbf{Required Attributes (minimum)} \\
\midrule
\endhead
User & Id, name, email (unique), password hash, role (\texttt{admin}/\texttt{teacher}/\texttt{student}) \\
Student & Id, linked user, admission number, guardian contact, enrollment date \\
Course & Id, code, name, assigned teacher, credits, fee amount, term \\
Enrollment & Id, student, course, status, unique per (student, course) \\
Grade & Id, enrollment, assessment type, score (0--100), recorded by \\
Invoice & Id, student, enrollment (unique -- one invoice per enrollment), amount (fee snapshot), method, status, issued/paid timestamps \\
Payroll Run & Id, period start/end, status, created by (planned feature) \\
Payslip & Id, payroll run, staff, base salary, deductions, net pay, status (planned feature) \\
\bottomrule
\end{longtable}

\section{Data Retention and Integrity}
\begin{itemize}[itemsep=2pt]
    \item An enrollment's status transitions (\texttt{active} $\rightarrow$ \texttt{completed}/\texttt{dropped}) shall be retained rather than deleted, preserving academic history.
    \item An invoice's \texttt{amount} shall remain fixed at the value snapshotted from the course fee at enrollment time, even if the course's fee later changes.
    \item Backups of the database shall be taken and verified by restore, per the Week 2 database-management build tasks.
\end{itemize}

% ================================================================
\chapter{Traceability and Closing Notes}

\section{Traceability to the Software Design Document}
Each functional requirement in Chapter~3 traces to a corresponding element
of the \sysname{} SDD: authentication requirements map to the SDD's Auth
Service design and sequence diagram (register/login); Academic and Finance
requirements map to the microservice architecture diagram and the
enrollment/invoice sequence diagram; Gateway requirements map to the
Gateway component and routing table; non-functional requirements map to the
SDD's Security, Performance, and Reliability sections.

\section{Traceability to the Exam Mark Scheme}
\begin{center}
\begin{tabular}{@{}p{3.6cm}p{4.4cm}p{6.4cm}@{}}
\toprule
\textbf{Exam Week} & \textbf{Marks} & \textbf{Requirements Covered} \\
\midrule
Week 1 -- Environment \& Auth & 5 & FR-AUTH-1 -- FR-AUTH-3, NFR-SEC-2 \\
Week 2 -- Database & 5 & Chapter 5 (Data Requirements) \\
Week 3 -- Microservices \& Gateway & 5 & FR-GW-1 -- FR-GW-4, FR-AUTH-6 -- FR-AUTH-7 \\
Week 4 -- Messaging, Security \& Performance & 5 & FR-ACAD-8, FR-FIN-1 -- FR-FIN-2, NFR-SEC-1 -- NFR-SEC-4, NFR-PERF-1 -- NFR-PERF-2, NFR-REL-1 -- NFR-REL-4 \\
Week 5 -- Module Feature, Testing \& CI/CD & 5 & FR-HR-2 -- FR-HR-4 or FR-FIN-5 (module feature), automated tests, CI pipeline (\textit{planned}) \\
Week 6 -- Integration \& Documentation & 5 & Full \texttt{docker-compose} bring-up, this SRS/SDD pair, all UML diagrams (\textit{planned}) \\
\bottomrule
\end{tabular}
\end{center}

\section{Open Items for the Group}
\begin{itemize}[itemsep=2pt]
    \item Decide which single feature the group will complete for Week 5's ``one complete feature for your assigned module'' task: payroll computation (HR), invoice payment/reconciliation (Finance), or a grade transcript (Academic).
    \item Set up the CI pipeline (GitHub Actions) and a first batch of automated tests for the chosen Week 5 feature.
    \item Decide, before Week 6, whether payment-gateway integration is added or explicitly documented as out of scope in the final submission.
\end{itemize}

\end{document}
